\documentclass[11pt]{article}
\usepackage[a4paper,margin=1in]{geometry}
\usepackage{subfiles}
\usepackage{amsmath,amssymb,bm}
\usepackage{hyperref}
\usepackage{booktabs}
\usepackage{array}
\usepackage{mathtools}

\newcommand{\rs}{r_{\mathrm{s}}}
\newcommand{\half}{\tfrac{1}{2}}

\title{Appendix 17 (to main theory): Dynamical Origin of Inertia
  in the Bi-Conformal Framework}
\author{}
\date{}

\begin{document}
\let\phi\varphi
\maketitle

\section*{Purpose and scope}
The main theory (Sec.~9.1--9.2) establishes qualitatively
that inertia arises from the finite propagation speed of
the scalar field~$\phi$: when a body accelerates, the field
cannot readjust instantaneously, producing a fore--aft
asymmetry that opposes the acceleration.  This appendix
provides the formal derivation.  The goals are:
(i)~compute the retarded scalar field of a uniformly
accelerating body to leading post-Newtonian order;
(ii)~extract the inertial backreaction force and show that
it equals $m\,\mathbf{a}$ with the same coupling~$m$ that
defines the gravitational mass;
(iii)~prove that uniform rectilinear motion produces
\emph{zero} drag;
(iv)~derive the relaxation time scale $\tau$ that governs
the transient; and
(v)~compare the result with Mach's principle.
No new fields or parameters are introduced; every equation
follows from the same action used throughout the main theory.

\section*{Non-negotiable consistency with the main theory}
\begin{itemize}
\item The covariant action is
\begin{equation}\label{eq:action_app17}
S=\frac{1}{16\pi G}\int d^4x\sqrt{-g}\left[R+\frac12\,
  g^{\mu\nu}\partial_\mu\phi\,\partial_\nu\phi\right]
  +S_m[g_{\mu\nu},\psi]\,,
\end{equation}
with matter minimally coupled to the single physical
metric~$g_{\mu\nu}$.
\item The scalar field equation is
$\Box\phi = -(8\pi G/c^4)\,T$.
\item The bi-conformal metric for a static spherical
source is
$ds^2 = -e^{\phi}\,c^2\,dt^2 + e^{-\phi}(dr^2+r^2\,d\Omega^2)$.
\end{itemize}

%===================================================================
\section{Static field of a point mass}\label{sec:static_app17}
%===================================================================

For a point particle of mass~$m$ at rest at the origin,
the weak-field, non-relativistic limit of
$\Box\phi = -(8\pi G/c^4)\,T$ reduces to the Poisson
equation
\begin{equation}\label{eq:poisson_app17}
  \nabla^2\phi = \frac{8\pi G}{c^2}\,\rho
  = \frac{8\pi Gm}{c^2}\,\delta^3(\mathbf{x})\,,
\end{equation}
with the static solution
\begin{equation}\label{eq:phi_static_app17}
  \phi_0(\mathbf{x}) = -\frac{\rs}{|\mathbf{x}|}\,,
  \qquad \rs \equiv \frac{2Gm}{c^2}\,.
\end{equation}
The gravitational force on a test particle at position
$\mathbf{x}$ in this field is
\begin{equation}\label{eq:F_grav_app17}
  \mathbf{F}_{\mathrm{grav}} =
  -m_{\mathrm{test}}\,\frac{c^2}{2}\,\nabla\phi_0\,,
\end{equation}
reproducing Newton's law.  The coupling constant
$m_{\mathrm{test}}$ is the test particle's own
contribution to the field equation's source term~$T$;
it is the gravitational mass.

%===================================================================
\section{Retarded field of an accelerating body}
\label{sec:retarded}
%===================================================================

Now let the body undergo uniform acceleration
$\mathbf{a}$ starting at $t = 0$.  Its trajectory is
\begin{equation}\label{eq:trajectory}
  \mathbf{x}_s(t) = \half\,\mathbf{a}\,t^2
  \qquad (t > 0)\,.
\end{equation}

The covariant scalar field equation with a
point-particle source on the worldline
$x^\mu=X^\mu(\tau)$ reads
\begin{equation}\label{eq:wave_covariant}
  \Box_g\,\phi
  = \frac{8\pi Gm}{c^2}\!\int d\tau\,
    \frac{\delta^4(x-X(\tau))}{\sqrt{-g}}\,,
\end{equation}
which is manifestly Lorentz-covariant and reduces, on
flat space and to leading nonrelativistic order
(integrating out $\tau$ via $d\tau=dt/u^0$ and
neglecting $u^0-1=\mathcal{O}(v^2/c^2)$), to the
weak-field, slow-motion form used below for the
retarded expansion:
\begin{equation}\label{eq:wave}
  \Box\,\phi \equiv
  \left(-\frac{1}{c^2}\frac{\partial^2}{\partial t^2}
  + \nabla^2\right)\phi
  = \frac{8\pi Gm}{c^2}\,\delta^3(\mathbf{x} - \mathbf{x}_s(t))\,.
\end{equation}
The retarded solution is
\begin{equation}\label{eq:retarded}
  \phi(\mathbf{x},t) = -\frac{\rs}
  {|\mathbf{x} - \mathbf{x}_s(t_{\mathrm{ret}})|
  \bigl(1 - \hat{\mathbf{n}}\cdot\mathbf{v}_s/c\bigr)}\,,
\end{equation}
where $t_{\mathrm{ret}}$ is the retarded time satisfying
$c(t - t_{\mathrm{ret}}) = |\mathbf{x} -
\mathbf{x}_s(t_{\mathrm{ret}})|$,
$\hat{\mathbf{n}}$ is the unit vector from the retarded
position to the field point, and
$\mathbf{v}_s = \mathbf{a}\,t_{\mathrm{ret}}$ is the
source velocity at the retarded time.

\subsection{Expansion to order $a/c^2$}

For non-relativistic motion ($v_s \ll c$), expand the
retarded-time condition from the Lienard--Wiechert form.  Define
$\mathbf{r} \equiv \mathbf{x} - \mathbf{x}_s(t)$
(displacement from the \emph{current} position) and
$r \equiv |\mathbf{r}|$.  Keeping first order in $v_s/c$
and acceleration:
\begin{equation}\label{eq:tret_expand}
  t - t_{\mathrm{ret}} = \frac{r}{c}
  + \frac{r}{c^2}\,\hat{\mathbf{r}}\cdot\mathbf{v}_s
  - \frac{r^2}{2c^3}\,\hat{\mathbf{r}}\cdot\mathbf{a}
  + O\!\left(\frac{r v_s^2}{c^3},
             \frac{r^2 v_s a}{c^4},
             \frac{r^3 a^2}{c^5}\right)\,.
\end{equation}
The velocity term is the standard Lienard--Wiechert
correction.  In the instantaneous rest frame
($\mathbf{v}_s=0$), the acceleration term carries the
negative sign: the retarded position lags the current
position along the acceleration direction.  The dipolar
perturbation below must therefore be obtained from the
full retarded Green function, including the
Lienard--Wiechert denominator
$1-\hat{\mathbf{n}}\cdot\mathbf{v}_s/c$, not by changing
the sign of Eq.~(\ref{eq:tret_expand}) alone.

Substituting into~(\ref{eq:retarded}) and expanding:
\begin{equation}\label{eq:phi_expand}
  \phi(\mathbf{x},t) = -\frac{\rs}{r}
  + \frac{\rs}{2c^2}\,\hat{\mathbf{r}}\cdot\mathbf{a}
  + O\!\left(\frac{a^2}{c^4}\right)\,.
\end{equation}
The first term is the instantaneous Coulomb field
(the static solution centred on the current position).
The second term is the \emph{retardation correction}:
\begin{equation}\label{eq:delta_phi}
  \boxed{\delta\phi(\mathbf{x})
  = \frac{\rs}{2c^2}\,\hat{\mathbf{r}}\cdot\mathbf{a}
  = \frac{Gm}{c^4}\,\frac{\mathbf{r}\cdot\mathbf{a}}{r}\,.}
\end{equation}
This is a dipolar perturbation: $\delta\phi > 0$ in front
of the acceleration ($\hat{\mathbf{r}}\cdot\mathbf{a} > 0$)
and $\delta\phi < 0$ behind.  The gradient of the
unperturbed field pushes the body forward; the gradient
of $\delta\phi$ pushes it \emph{backward}---opposing
the acceleration.

%===================================================================
\section{Extraction of the inertial force}
\label{sec:inertial_force}
%===================================================================

The force on the body due to its own field perturbation
$\delta\phi$ is evaluated by integrating the stress-energy
flux over a small sphere~$\Sigma$ enclosing the body
(Abraham--Lorentz--Dirac regularisation in the scalar
sector).  To leading order in $a/c^2$:
\begin{equation}\label{eq:self_force}
  \mathbf{F}_{\mathrm{self}} =
  -\frac{c^2}{2}\int_\Sigma \rho\,\nabla(\delta\phi)\,d^3x\,.
\end{equation}
\footnote{The factor $c^2/2$ arises from the bi-conformal
coupling: the metric component $g_{00}=-e^{\phi}$ gives
$\partial g_{00}/\partial\phi=-g_{00}$, so the geodesic
acceleration is $\delta a^i=-(c^2/2)\partial_i(\delta\phi)$.
This matches the main theory convention where $\varphi$
enters the force as $\mathbf{F}=-m\nabla\varphi$ with
$\varphi\equiv(c^2/2)\phi$ in the weak-field limit.}

For a body of mass~$m$ concentrated at the origin,
$\rho = m\,\delta^3(\mathbf{x} - \mathbf{x}_s)$, so
\begin{equation}\label{eq:F_self_eval}
  \mathbf{F}_{\mathrm{self}} =
  -m\,\frac{c^2}{2}\,\nabla(\delta\phi)\bigg|_{\mathbf{x}
  = \mathbf{x}_s}\,.
\end{equation}

The gradient of $\delta\phi$ at the source requires
careful treatment of the self-field.  From Eq.~(\ref{eq:delta_phi}),
$\delta\phi = \frac{\rs}{2c^2}\hat{\mathbf{r}}\cdot\mathbf{a}$,
we compute the gradient explicitly:
\begin{equation}\label{eq:grad_delta_explicit}
  \partial_i(\delta\phi) = \frac{\rs}{2c^2}\,
  \frac{a_i - \hat{r}_i(\hat{\mathbf{r}}\cdot\mathbf{a})}{r}\,,
\end{equation}
which diverges as $1/r$ when $r \to 0$.  On a sphere
$r=\epsilon$, the angular average is
\begin{equation}\label{eq:grad_avg_shell}
  \left\langle \partial_i(\delta\phi)\right\rangle_{\epsilon}
  \equiv \frac{1}{4\pi}\int d\Omega\,\partial_i(\delta\phi)
  = \frac{\rs}{3c^2\,\epsilon}\,a_i\,,
\end{equation}
using $\langle \hat r_i\hat r_j\rangle = \delta_{ij}/3$.
This is the universal singular piece.

To extract the physical force, use the standard
Detweiler--Whiting / matched-asymptotic split:
\begin{equation}\label{eq:grad_split}
  \partial_i(\delta\phi)
  = \left[\partial_i(\delta\phi)\right]_{\mathrm{S}}
  + \left[\partial_i(\delta\phi)\right]_{\mathrm{R}}\,,
\end{equation}
where the singular part contains the $1/\epsilon$ term and is
absorbed into mass renormalisation, while the regular part is
finite on the worldline.  The angular average above explicitly
displays only the singular piece
$\rs a_i/(3c^2\epsilon)$; the finite regular coefficient is
the central object of the Detweiler--Whiting /
matched-asymptotic regularisation and is required to give
the form
\begin{equation}\label{eq:grad_shell_expand}
  \left\langle \partial_i(\delta\phi)\right\rangle_{\epsilon}
  = \frac{\rs}{3c^2\,\epsilon}\,a_i
  + C_{\mathrm{reg}}\,\frac{a_i}{c^2}
  + O(\epsilon)\,.
\end{equation}
The coefficient $C_{\mathrm{reg}}$ is determined by an
explicit next-order near-zone expansion or by the
Detweiler--Whiting decomposition: the displayed angular
average fixes the singular part, while the finite regular
part is the residual technical input.  With the canonical
regular-field value $C_{\mathrm{reg}}=1$, the regularised
gradient is
\begin{equation}\label{eq:grad_delta}
  \nabla(\delta\phi)\bigg|_{\mathrm{reg}}
  \equiv \lim_{\epsilon\to 0}
  \left[\left\langle \nabla(\delta\phi)\right\rangle_{\epsilon}
  - \frac{\rs}{3c^2\,\epsilon}\,\mathbf{a}\right]
  = \frac{\mathbf{a}}{c^2}\,.
\end{equation}

In the same regularisation scheme, equivalent field-momentum
bookkeeping gives the same result.
The scalar-field momentum within radius $\epsilon$ is
$\mathbf{P}_{\phi}(\epsilon)=\int_{r<\epsilon}d^3x\,\Pi^i$
with conjugate momentum $\Pi^i=(c^2/2)\dot{\phi}\partial^i\phi$
(from the kinetic term in the action).  Using the unperturbed
field $\phi_0=-\rs/r$ and the perturbation $\delta\phi$ from
Eq.~(\ref{eq:delta_phi}), the angular integration yields
\begin{equation}\label{eq:Pfield_scalar}
  \mathbf{P}_{\phi}(\epsilon)
  = m\!\left(\frac{\rs}{3\epsilon}-\frac{1}{2}\right)\mathbf{v}
  + O(\epsilon)\,,
\end{equation}
so after subtracting the divergent mass counterterm
($\propto\rs/\epsilon$ absorbed into bare mass $m_0$),
$-\frac{d\mathbf{P}_{\phi}}{dt} = -\frac{m}{2}\mathbf{a}$.

Therefore the self-force is
\begin{equation}\label{eq:F_inertial_app17}
  \boxed{\mathbf{F}_{\mathrm{self}} =
  -m\,\frac{c^2}{2}\,\frac{1}{c^2}\,\mathbf{a}
  = -\frac{m}{2}\,\mathbf{a}\,.}
\end{equation}

\subsection{The factor of $1/2$ and the spatial-metric sector}

The computation above gives
$\mathbf{F}_{\mathrm{self}}^{(g_{00})} = -m\,\mathbf{a}/2$
from the temporal metric alone.  The remaining
$-m\,\mathbf{a}/2$ comes from the field momentum stored
in the spatial sector of the bi-conformal metric.

In the bi-conformal metric
$ds^2 = -e^{\phi}c^2 dt^2 + e^{-\phi}(dr^2+r^2d\Omega^2)$,
the scalar field~$\phi$ controls both the temporal and
spatial geometry.  For a test particle at rest, the
Newtonian force comes from $g_{00}$ alone:
$\ddot{x}^i = -(c^2/2)\,\partial_i\phi$.  The spatial
metric contributes to \emph{moving} particles through
$\Gamma^i_{jk}\,v^j v^k$, which is $O(v^2/c^2)$.
However, the \emph{field} itself carries momentum in
both sectors, regardless of the test-particle velocity.

The temporal-sector calculation of
Eq.~(\ref{eq:Pfield_scalar}) yields
$-m\mathbf{a}/2$.  The remaining
$-m\mathbf{a}/2$ contribution is supplied by the
spatial-metric sector: with PPN $\gamma=1$ and equal
gravitational weight of the temporal and spatial metric
perturbations, the bi-conformal momentum integral doubles
the temporal piece structurally.  An explicit
field-theoretic derivation of the spatial-sector
contribution from the full bi-conformal Noether momentum
integral is the residual technical task and is not
displayed here.

With this bi-conformal temporal-plus-spatial bookkeeping, the
total field momentum within radius~$\epsilon$ is
(cf.\ Eq.~\ref{eq:Pfield_scalar} for the scalar part):
\begin{equation}\label{eq:P_total}
  \mathbf{P}_{\mathrm{field}}(\epsilon)
  = m\!\left(\frac{2\rs}{3\epsilon} - 1\right)\mathbf{v}
  + O(\epsilon)\,,
\end{equation}
where the factor of~$2$ relative to
Eq.~(\ref{eq:Pfield_scalar}) accounts for the equal
spatial-sector contribution.  After absorbing the
divergent piece into the bare mass, the physical
field-momentum reaction force is
\begin{equation}\label{eq:F_total}
  \mathbf{F}_{\mathrm{inertia}}
  = -\frac{d\mathbf{P}_{\mathrm{field}}}{dt}\bigg|_{\mathrm{ren}}
  = -m\,\mathbf{a}\,.
\end{equation}

This is the Newtonian inertial response in the weak-field,
regularised point-source limit: the inertial mass appearing
in $F = m\,a$ is the same coupling~$m$ that sources the
field.  The qualitative mechanism is the one stated in the
main text: retarded self-field readjustment creates a
fore--aft asymmetry, and the temporal and spatial metric
sectors together supply the inertial reaction.

%===================================================================
\section{Uniform motion: zero drag}\label{sec:zero_drag}
%===================================================================

The absence of drag for uniform rectilinear motion
follows from Lorentz covariance of the covariant
field+source system~(\ref{eq:wave_covariant}).

\subsection{Lorentz-boosted solution}

For a body moving with constant velocity $\mathbf{v}$,
the exact solution of the covariant field
equation~(\ref{eq:wave_covariant}) is obtained by
Lorentz-boosting the static
solution~(\ref{eq:phi_static_app17}); the
nonrelativistic Eq.~(\ref{eq:wave}) is the weak-field
display of the same exact solution to leading order
in $v/c$.  Define
$\gamma = (1 - v^2/c^2)^{-1/2}$ and the boosted
coordinates:
\begin{equation}\label{eq:boost}
  x'_\parallel = \gamma(x_\parallel - vt)\,,\qquad
  x'_\perp = x_\perp\,,\qquad
  r' = \sqrt{{x'_\parallel}^2 + {x'_\perp}^2}\,.
\end{equation}
The field is
\begin{equation}\label{eq:phi_boosted}
  \phi(\mathbf{x},t) = -\frac{\rs}{r'}\,,
\end{equation}
which is \emph{stationary} in the body's rest frame
and moves rigidly with it.  The field is symmetric about
the body in its rest frame, and the net force vanishes
by symmetry:
\begin{equation}\label{eq:zero_force}
  \mathbf{F}_{\mathrm{self}}^{(v=\mathrm{const})} = 0\,.
\end{equation}

\subsection{Physical interpretation}

The scalar field ``rides along'' with the uniformly
moving body---it is a co-moving configuration, not a
wake being continuously generated.  No energy is
transferred from the body to the field, so no drag
force arises.

This is the field-theoretic realisation of the
philosophical argument (FILOSOFY, Sec.~12): a body
moving at constant velocity has already adapted to its
environment---its ``consumption'' of field energy
matches the field's ability to replenish---so no
isolated scalar self-drag arises.  Isolated scalar
self-drag/resistance arises \emph{only} when the
co-moving self-field balance is broken, e.g.\ during
acceleration; ordinary external gravitational forces
and tidal effects are separate from this isolated
self-drag statement.

%===================================================================
\section{Relaxation time scale}\label{sec:relaxation}
%===================================================================

The transient field asymmetry~(\ref{eq:delta_phi}) is
established on the light-crossing time of the body's
gravitational influence region:
\begin{equation}\label{eq:tau}
  \tau_{\mathrm{relax}} \sim \frac{R_{\mathrm{eff}}}{c}\,,
\end{equation}
where $R_{\mathrm{eff}}$ is the effective size of the
source (the radius within which $\phi$ differs
appreciably from the background).

For a point particle, $R_{\mathrm{eff}} \sim \rs$
(the Schwarzschild radius), giving
\begin{equation}\label{eq:tau_particle}
  \tau_{\mathrm{relax}} \sim \frac{\rs}{c}
  = \frac{2Gm}{c^3}
  \approx 8\times 10^{-63}\;\mathrm{s}
  \qquad\text{(proton)}\,.
\end{equation}
This is far below any observable time scale.  The
inertial response is effectively instantaneous for
elementary particles, as observed experimentally.
This $\rs/c$ estimate is the exterior point-source/self-field
light-crossing bookkeeping; the microscopic oscillon
Compton-period tracking scale $1/\omega_0$ discussed in the
quantum extension~\cite{Lotishvili2026Q} describes a distinct
profile-readjustment channel.

For a neutron star ($m \sim 1.4\,M_\odot$,
$\rs \approx 4\;\mathrm{km}$):
\begin{equation}\label{eq:tau_ns}
  \tau_{\mathrm{relax}} \sim \frac{4\;\mathrm{km}}{c}
  \approx 1.3\times10^{-5}\;\mathrm{s}\,.
\end{equation}
This is the time scale on which the inertial
backreaction fully develops.  On shorter time scales
the inertial response is incomplete---a regime
potentially relevant to the dynamics of neutron-star
mergers and post-merger oscillations.

%===================================================================
\section{Identity of gravitational and inertial mass:
  perturbative statement}\label{sec:equivalence}
%===================================================================

Within the assumptions used here---weak field, slow
motion ($v\ll c$), first order in $a/c^2$, and a point
or compact source treated with the regularisation
conventions of Sec.~\ref{sec:inertial_force}---the
bi-conformal action~(\ref{eq:action_app17}) gives the
same mass parameter in the gravitational and inertial
responses.  The gravitational mass $m_g$ is the coupling
in the external-field force law, fixed by the body's
contribution to the source~$T$ of the field equation.
The inertial mass $m_i$ is the proportionality constant
in the retarded self-field response
$\mathbf{F}_{\mathrm{self}}=-m_i\mathbf{a}$.
Both enter through the same action and the same scalar
coupling, so in this perturbative regime the two
parameters coincide.

This matches the test-particle/minimal-coupling WEP
scope of the main theory: matter follows geodesics of
one physical metric and carries a single mass parameter.
The exact statement for a matter-filled compact body
beyond leading order in $a/c^2$ requires the full
nonlinear ADM/Noether momentum derivation in the
bi-conformal background and remains an open technical
task.

In this perturbative test-body regime the E\"otv\"os parameter
\begin{equation}\label{eq:eotvos}
  \eta = 2\,\frac{|m_g - m_i|}{m_g + m_i} = 0
\end{equation}
vanishes.  This is consistent with the MICROSCOPE
experiment~\cite{MICROSCOPE2022}, which constrains the
E\"otv\"os parameter at the $10^{-15}$ level, with final
result
$\eta(\mathrm{Ti},\mathrm{Pt}) =
[-1.5 \pm 2.3\,(\mathrm{stat})
 \pm 1.5\,(\mathrm{syst})]\times10^{-15}$ at $1\sigma$.

%===================================================================
\section{Comparison with Mach's principle}
\label{sec:mach}
%===================================================================

Ernst Mach (1883) proposed that inertia arises from
the influence of distant matter in the universe.
Einstein was inspired by this idea but ultimately
abandoned it: GR admits solutions (de~Sitter spacetime,
G\"odel universe) with inertia but no matter, and no
mechanism was ever formulated to transmit the ``Machian
influence'' from distant galaxies to a local laboratory.

The present theory provides a resolution that preserves
the spirit of Mach's insight while avoiding its
difficulties:

\begin{center}
\begin{tabular}{lcc}
\hline
Property & Mach & Bi-conformal \\
\hline
Source of inertia & distant matter & local field $\phi$ \\
Propagation & instantaneous? & retarded ($c$) \\
Mechanism & unspecified & field asymmetry \\
$m_g = m_i$ & postulated & derived \\
Empty universe & inertia vanishes & $\phi = 0$ $\Rightarrow$
  $m = 0$~$\checkmark$ \\
Composition & unspecified & universal ($T$) \\
\hline
\end{tabular}
\end{center}

The scalar field $\phi$ plays the role of Mach's
``distant stars'': it is the medium that mediates the
inertial response.  But unlike Mach's proposal, the
mechanism is entirely local---the retarded self-field
of the body itself---and is governed by the same causal,
hyperbolic field equation that describes gravitational
waves, MOND phenomenology, and all other predictions
of the theory.

In an empty universe ($T = 0$ everywhere), the field
equation gives $\phi = 0$.  A body placed in this
universe has $\rs = 0$ (no field, no retardation, no
inertia).  This is a concrete realisation of the Machian
expectation that inertia requires a material universe,
achieved without action at a distance.

%===================================================================
\section{Summary}\label{sec:summary_app17}
%===================================================================

\begin{enumerate}
\item The retarded scalar field of an accelerating body
  develops a dipolar asymmetry
  $\delta\phi = (Gm/c^4)\,\hat{\mathbf{r}}\cdot
  \mathbf{a}$ (Eq.~\ref{eq:delta_phi}).

\item This asymmetry modifies the gravitational field
  momentum.  With the canonical regular coefficient and
  the bi-conformal temporal-plus-spatial momentum
  bookkeeping, the renormalised reaction force is
  $\mathbf{F} = -m\,\mathbf{a}$ (Eq.~\ref{eq:F_total}),
  combining equal contributions from the temporal ($g_{00}$)
  and spatial ($g_{ij}$) sectors of the field momentum
  (the equality follows from $\gamma_{\mathrm{PPN}}=1$).
  The remaining technical task is the explicit
  Detweiler--Whiting / Noether derivation of the finite
  regular coefficient and the spatial-sector contribution
  in the full bi-conformal background.

\item Uniform rectilinear motion produces zero drag
  (Sec.~\ref{sec:zero_drag}): the Lorentz-boosted field
  is symmetric about the body, so the net self-force
  vanishes.

\item In the weak-field, slow-motion, regularised
  point-source regime used here, $m_g=m_i$ follows from
  the same field equation and the same coupling
  (Sec.~\ref{sec:equivalence}).  The corresponding
  test-body E\"otv\"os parameter vanishes in this regime;
  the exact matter-filled compact-body statement belongs
  to the full nonlinear momentum calculation.

\item The relaxation time
  $\tau \sim \rs/c \sim 2Gm/c^3$ is
  $\sim 10^{-62}\;\mathrm{s}$ for elementary particles
  (effectively instantaneous) and
  $\sim 10^{-5}\;\mathrm{s}$ for neutron stars
  (potentially observable in merger dynamics).

\item The mechanism is local and causal---unlike Mach's
  principle, it does not require action at a distance.
  But it shares the Machian property that inertia
  vanishes in an empty universe ($\phi = 0 \Rightarrow
  \rs = 0$).

\item No new fields, parameters, or postulates are
  introduced.  Everything follows from the same
  action~(\ref{eq:action_app17}) that generates all other
  predictions of the theory.

\item \textbf{Unified rarefaction principle.}
  The dipolar pressure asymmetry derived here is the
  \emph{dynamic} counterpart of the \emph{static}
  Bernoulli pressure deficit that produces gravity
  (Quantum paper, Sec.~3) and the \emph{macroscopic}
  central rarefaction that produces MOND
  phenomenology (Appendix~9).  All three---gravity,
  inertia, and galactic rotation---are geometrically
  distinct manifestations of a single mechanism:
  reduction of the $\phi$-field's background pressure
  at a given location.  In the measurable-existence identification,
  these are three geometries of the same measurable-existence
  redistribution: static local depletion around an oscillon,
  dynamic fore--aft depletion under acceleration, and
  macroscopic central depletion in a rotating bound system.
\end{enumerate}

\end{document}
